\subsection{Axiom Verification}\label{subsec:axioms}

The eight dm3 axioms of~\cite{Paper1} are verified explicitly for
the toy model on $M=\mathbb{R}^2_{>0}\times\mathbb{R}$ with
equations~\eqref{eq:rdot}--\eqref{eq:zdot}. Each axiom is checked
with explicit derivations and cross-references to the global results
of Sections~\ref{sec:invariants}--\ref{sec:measures}.

\paragraph{Axiom~1 (Orbit identity).}
The four phase points $p_I=(1,0)$, $p_{VI}=(1,\pi/2)$,
$p_{III}=(1,\pi)$, $p_{VII}=(1,3\pi/2)$ lie on $\Orb=\{r=1\}$
and are cycled by $\phi^{\pi/2}$ (Proposition~\ref{prop:axioms}).

\paragraph{Axiom~2 (Generative closure).}
On $\Orb$: $\dot r = 1\cdot(1-1) + 0 = 0$, so $r$ remains $1$
under the flow. $\Orb$ is compact, connected, and invariant.

\paragraph{Axiom~3 (Structural primacy).}
$V=(r-1)^2$. Computing:
\[
  \dot V = 2(r-1)\dot r
  = 2(r-1)\bigl[r(1-r^2) + 2(r-1)e^{-z}\bigr]
  = -2(r-1)^2\bigl[r(1+r) - 2e^{-z}\bigr].
\]
For $r$ near $1$ and $z>0$: $r(1+r)\approx 2 > 2e^{-z}$, so
$\dot V < 0$ for $r\ne 1$. Stability is encoded in the vector
field and $V$, not in external constraints (Axiom~7 holds
simultaneously).

\paragraph{Axiom~4 (Perturbation response).}
Near $\Orb$ with $\rho=r-1$ small, expand to first order in $\rho$:
\[
  \dot\rho = -2(1-e^{-z})\rho + O(\rho^2).
\]

\begin{lemma}[Transverse stability]\label{lem:transverse}
For $z>0$, the transverse eigenvalue is
$\lambda(z) = -2(1-e^{-z}) < 0$.
The orbit $\Orb$ is exponentially stable with rate
$c(z) = 2(1-e^{-z})$, which increases to $2$ as $z\to\infty$.
\end{lemma}

\begin{proof}
The linearization $\dot\rho = \lambda(z)\rho$ has solution
$\rho(t) = \rho_0\,e^{\lambda(z)t}\to 0$ for $\lambda(z)<0$,
i.e., for $z>0$.
\end{proof}

\begin{remark}[The embodiment threshold in the linearization]
At $z=0$: $\lambda(0) = 0$. The orbit is \emph{neutrally stable}
at the moment of first contact. Stability emerges only as $z$
accumulates, i.e., as the system traverses $\Orb$ and action
builds up. This is the precise dynamical content of the embodiment
threshold: the orbit becomes attracting only after it has been
traversed. The Lyapunov condition $\dot V\le -cV$ holds with
$c = c(z) = 2(1-e^{-z})$, yielding embodiment threshold
$\tau = \sqrt{c/\kappa} = \sqrt{4/1} = 2$ at the asymptotic
rate $c=4$ (Axiom~5 below).
\end{remark}

\paragraph{Axiom~5 (Noise-as-fuel).}
For the It\^o SDE $dr = \dot r\,dt + \sigma\,dW_t$, the generator
applied to $V=(r-1)^2$ gives
\[
  \mathcal{L}V = \dot V + \tfrac{\sigma^2}{2}\cdot 2
  = -4(r-1)^2(1-e^{-z}) + \sigma^2.
\]
Near $\Orb$ with $z>0$: $\mathcal{L}V \le -4V(1-e^{-z}) + \sigma^2$.
Setting $c=4(1-e^{-z})>0$ and $\kappa=1$:
\[
  \tau = \sqrt{c/\kappa} = 2\sqrt{1-e^{-z}}\;\nearrow\; 2
  \quad\text{as }z\to\infty.
\]
The asymptotic embodiment threshold is $\tau=2$.

\begin{corollary}[Noise tolerance under unification]
By Theorem~B of~\cite{Paper1}, unification of two toy model
systems gives $\tau_{12}\le\min(\tau_1,\tau_2)=2$.
\end{corollary}

\paragraph{Axiom~6 (Non-collapse).}
$\dot\theta=1\ne 0$ on $\Orb$; minimal period $T^*=2\pi>0$.

\paragraph{Axiom~7 (Non-coercion).}
The vector field $(r(1-r^2)+2(r-1)e^{-z},\,1,\,r^2-2(r-1)^2e^{-z})$
is $C^\infty$ on $M$. No discontinuous projection or hard boundary
enforcement appears.

\paragraph{Axiom~8 (Hyperbolicity).}
From Lemma~\ref{lem:transverse}: for $z>0$, the transverse
Floquet exponent is $\mu_{\max}=-2<0$. All non-trivial Floquet
multipliers satisfy $\abs{\mu_j}<1$.

\medskip

The canonical invariant triple is therefore
\[
  \iota(\dm_0) = (T^*,\,\mu_{\max},\,\tau)
  = (2\pi,\,-2,\,2),
\]
and the structural stability radius is
\[
  \varepsilon_0
  = \frac{\abs{\mu_{\max}}}{2(1+\sup_\Orb\norm{\Hess V}_\infty)}
  = \frac{2}{2(1+2)}
  = \frac{1}{3}.
\]

All eight axioms are satisfied. The toy model is a complete, explicit
instantiation of the dm3 framework of~\cite{Paper1}. \qed
